\section{Background and related work}
\label{sec:related}

\subsection{Pareto efficiency and \emph{post hoc} selection}
For a vector objective $F(x)=(f_1(x),\dots,f_m(x))$ (minimisation), $x$
Pareto-dominates $y$ if $f_i(x)\le f_i(y)$ for all $i$ with at least one strict
inequality; the efficient set is the non-dominated set. Because it is typically
large, a single recommendation is obtained by \emph{post hoc} selection. The
weighted sum, TOPSIS \citep{hwang1981}, compromise programming
\citep{zeleny1982}, knee selection \citep{branke2008}, ELECTRE \citep{roy1968,roy1991},
PROMETHEE \citep{brans2016} and hypervolume selection \citep{zitzler1999} are the
standard tools; each requires weights, thresholds, a metric, or an already
generated front. LexUR does not eliminate modelling choices. It relocates them into a declared family of admissible monotone probes and, on finite candidate sets, computes the resulting recommendation exactly while returning a labelled regret certificate.

\subsection{Achievement scalarising functions and ordered aggregation}
Achievement scalarising functions (ASF) project a reference point onto the
efficient set through a parameterised metric interpolating between $L_1$ and
$L_\infty$ \citep{wierzbicki1980,miettinen2002}. Ordered weighted averaging
(OWA) \citep{yager1988} and its extensions (WOWA) \citep{torra1997}
connect ordered aggregation to lexicographic maximin. 

\begin{remark}[Recovery of standard scalarisations]
\label{rem:special}
The LexUR core recovers standard aggregation methods as exact special cases under specific, restricted probe families:
\begin{itemize}
\item \textbf{Weighted sum:} If $\Q=\{q_w(r)=w^\top r\}$, LexUR exactly minimises the weighted sum of normalised criteria.
\item \textbf{Chebyshev / ASF:} If $\Q=\{q_w(r)=\max_i w_i r_i\}$, LexUR exactly minimises the Chebyshev distance to the ideal point.
\item \textbf{Lexicographic minimax:} If $\Q=\{r_1,\dots,r_m\}$ (the singleton family), LexUR performs a lexicographic minimax optimisation over the normalised regrets.
\end{itemize}
\end{remark}

We do not claim
LexUR subsumes these methods with its full adaptive family---there it is a
\emph{different} aggregation paradigm. Its distinguishing features are that it
sorts \emph{normalised regrets} (dimensionless disappointments relative to each
probe's own attainable best and worst) over a \emph{declared family of probes}
and breaks ties through the whole sorted vector, and that it emits a certificate.

\subsection{Robust MCDA: SMAA, robust ordinal regression, preference robust
optimisation}
The work closest to LexUR lies in \emph{robust} multicriteria decision analysis and minimax regret optimisation \citep{kouvelis1997}.
\textbf{SMAA} \citep{lahdelma1998,tervonen2007} computes, by Monte-Carlo
integration over a distribution of admissible weights (and possibly criterion
values), acceptability indices describing how strongly each alternative is
supported. \textbf{Robust ordinal regression} \citep{greco2008,figueira2009}
reasons over the entire set of additive value functions compatible with a DM's
holistic pairwise judgements, producing necessary and possible preference
relations, sharing goals with inverse MCDA and the UTA family \citep{jacquetlagreze1982}. \textbf{Preference robust optimisation} \citep{armbruster2015,hu2015,delage2015,vayanos2020}
optimises the worst-case expected utility over an ambiguity set of utility
functions. LexUR shares the robust philosophy---perform well across a whole class of
admissible preferences rather than commit to one---but differs in three ways.
First, its admissible class is an explicit family of \emph{monotone probes}, not a
weight distribution or an elicited polyhedron. Second, its criterion is the
\emph{lexicographic minimax of normalised regret}, not an expected acceptability,
a necessary/possible relation, or a single worst-case utility. Third, its output
is an interpretable \emph{certificate} naming the binding coalitions. We note that
SMAA is inherently an \emph{exploratory} framework designed to visualise the robustness of alternatives across weight spaces, rather than a single-choice prescriptive rule. By forcing it to recommend the alternative with the highest first-rank acceptability index (ties broken by expected value), we create a comparable baseline for \S\ref{sec:exp}; however, we emphasise that LexUR's single-shot default recommendation is meant to \emph{complement} rather than compete with SMAA's rich exploratory indices. We likewise compare against a Savage minimax-regret (MMR)
recommender over a sampled linear weight set (\S\ref{sec:exp}). More broadly, LexUR
is complementary to data-driven preference disaggregation and inverse MCDA
\citep{doumpos2011}, which infer preference parameters from holistic judgements:
LexUR instead dispenses with a fitted preference model and reasons over a declared
class of questions, and could take an elicited probe family as input where such
data exist.

\subsection{Minimax regret and robust optimisation}
Regret-based choice originates with \citet{savage1951} and \citet{loomes1982};
minimax regret seeks alternatives minimising the largest shortfall from the best
attainable value. Robust optimisation \citep{bental2002} and stochastic
programming \citep{birge2011} address parameter uncertainty in the model rather
than preference uncertainty in the DM. LexUR extends minimax regret in three
directions: it is \emph{lexicographic} (after the worst regret, minimise the
second worst, and so on), \emph{probe-based} (regret is computed over rational
coalitions of criteria, not raw criteria), and \emph{direct on a given candidate
set} (computed without an additional scalarisation sweep, and front-free only in
the structured continuous cases of \S\ref{sec:direct}).

\subsection{Social choice and fairness}
Aggregating criteria resembles social choice, but treating criteria as voters
invites Condorcet cycles; LexUR avoids this by aggregating \emph{regrets over
declared questions}, not votes over criteria. For multi-stakeholder settings we
draw on the maximin principle of \citet{rawls1971} and the bargaining viewpoint
of \citet{nash1950}; we claim natural fairness properties (anonymity, monotonicity,
worst-stakeholder protection), not a uniqueness theorem.

\subsection{What is new relative to existing regret and robust-MCDA rules}
\label{sec:novelty}

While LexUR shares motivation with several robust approaches, its mathematical object and output differ fundamentally:
\begin{itemize}
\item \textbf{Versus minimax regret:} LexUR minimises a sorted vector of normalised probe regrets, not a single maximum regret over raw utilities or sampled weights.
\item \textbf{Versus ASF/Chebyshev:} ASF commits to one scalarising function or one reference structure; LexUR evaluates a declared family and certifies binding probes.
\item \textbf{Versus SMAA:} SMAA integrates over a preference distribution and reports acceptability; LexUR does not require a probability distribution and returns a leximax regret certificate.
\item \textbf{Versus robust ordinal regression (ROR):} ROR reasons over compatible value functions after elicited comparisons; LexUR can operate without holistic preference judgements.
\item \textbf{Versus preference robust optimisation (PRO):} PRO optimises worst-case utility over an ambiguity set; LexUR optimises lexicographically sorted normalised disappointments over interpretable monotone probes.
\item \textbf{Versus ordered weighted aggregation (OWA):} LexUR sorts regrets
across declared questions, not criterion values themselves. OWA specifies an
ordered weight vector, whereas LexUR specifies a probe family; both therefore
encode modelling choices in different forms. We include OWA utilities as a
held-out evaluation family, not as evidence that LexUR subsumes OWA.
\end{itemize}

The novelty is not minimax regret alone, but lexicographic minimax over normalised disappointments induced by a declared monotone probe family, with an auditable certificate and an adaptive redundancy-aware probe construction.

Table~\ref{tab:positioning} summarises where LexUR sits relative to representative MCDA and robust methods.

\begin{table}[t]
\centering\small
\caption{Comparison of LexUR with representative selection and robust-MCDA methods.}
\label{tab:positioning}
\begin{tabular}{p{2.2cm}p{3.2cm}p{1.3cm}p{2.6cm}p{3.7cm}}
\toprule
Method & Required Input & Pareto Comp. & Output / Explanation & Key Limitation \\
\midrule
Weighted sum & Weight vector & Yes ($w>0$) & Single choice (Low) & Misses non-supported points in non-convex fronts \\
TOPSIS & Weights, distance metric & Not strictly & Single choice (Low) & Rank reversal; distance to non-feasible ideal \\
CP & Weights, $L_p$ exponent & Yes ($p<\infty$) & Single choice (Low) & Non-linear scaling effects \\
ASF / Chebyshev & Reference point, scaling & Yes & Single choice (Low) & Requires an arbitrary reference point \\
MMR (sampled) & Linear weight distribution & Yes & Worst-case regret & Limited to linear aggregations \\
SMAA & Preference distribution & Usually & Acceptability indices & Exploratory; single-choice depends on post-processing \\
VIKOR & Weights, majority threshold & Not strictly & Single choice (Low) & Complex parameter interaction \\
Single-point HV & Reference point & Yes & Single choice (Low) & Box-volume score; sensitive to reference point \\
\textbf{LexUR (ours)} & Monotone probe family & Yes\textsuperscript{*} & Leximax certificate & Sensitive to normalisation bounds \\
\bottomrule
\multicolumn{5}{l}{\footnotesize \textsuperscript{*} Strict Pareto compatibility holds if probes separate criteria (Theorem~\ref{thm:pareto}).}
\end{tabular}
\end{table}
